\begin{ouverture}[calcul]{2}{Additionner et soustraire}{Partie 1 · Nombres entiers et calculs}
{Sommes · Différences · Estimation · Schémas en barres}
\objectif{Poser et effectuer additions et soustractions d’entiers.}
\objectif{Estimer un résultat avant de calculer.}
\objectif{Résoudre des problèmes en une, deux ou trois étapes.}
\prerequis{Numération jusqu’au million et tables d’addition.}
\end{ouverture}
\pagestyle{otMain}\essentiel\fontsize{11}{15}\selectfont
\phantomsection\label{cours-2-1}
\essSection{1}{Calculer des sommes et des différences}
\begin{definition}Le résultat d’une addition est une somme ; les nombres ajoutés sont les termes. Le résultat d’une soustraction est une différence. Dans une addition, l’ordre des termes ne change pas la somme : $47+385=385+47$.\end{definition}
\begin{savoirfaire}[Poser proprement]Pour poser une opération, aligner les chiffres rang par rang : unités sous unités, dizaines sous dizaines. Dans $5\,204-1\,368$, un rang trop petit s’échange : une dizaine contre dix unités. On obtient $3\,836$.\end{savoirfaire}
\begin{exemple}$6\,478+525$ : $8+5=13$, on écrit 3 et on retient une dizaine. En poursuivant rang par rang, la somme est $7\,003$. Contrôle : $7\,003-525=6\,478$.\end{exemple}
\begin{savoirfaire}[Estimer puis contrôler]Avant de calculer, arrondir chaque terme : $6\,478+525$ est proche de $6\,500+500=7\,000$. Une réponse éloignée de l’estimation doit être refaite. Pour ajouter 9, 19 ou 29 de tête : ajouter 10, 20 ou 30, puis retirer 1.\end{savoirfaire}
\clearpage
\phantomsection\label{cours-2-2}
\essSection{2}{Résoudre un problème pas à pas}
\begin{propriete}Résoudre un problème se fait en quatre étapes : comprendre l’histoire du problème, modéliser en choisissant les opérations, calculer, puis répondre par une phrase et vérifier que la réponse est vraisemblable.\end{propriete}
\label{fig-cours-2-2-cours-barres}\figureGeo{cours-barres}

\begin{exemple}« Une école compte 258 élèves ; 136 sont des filles. » Le schéma en barres montre un tout (258) et deux parties. Le nombre de garçons est la partie manquante : $258-136=122$.\end{exemple}
\begin{attention}Un mot ne décide pas de l’opération. « Léa a 12 billes de plus que Tom » peut demander une soustraction, si on connaît le nombre de billes de Léa et qu’on cherche celui de Tom. Relire l’histoire entière avant de calculer.\end{attention}

\clearpage\rubrique{Découverte}
\begin{blocLivre}{ex-2-1}\begin{exercice}[1]{Deux sommes posées}
Poser et calculer $6\,478+525$ et $38\,467+4\,850$.
\end{exercice}\end{blocLivre}
\begin{blocLivre}{ex-2-2}\begin{exercice}[1]{Deux différences posées}
Poser et calculer $5\,204-1\,368$ et $70\,000-2\,460$.
\end{exercice}\end{blocLivre}
\begin{blocLivre}{ex-2-3}\begin{exercice}[1]{Compléments utiles}
Compléter : $64+\ldots=100$, puis $380+\ldots=1\,000$.
\end{exercice}\end{blocLivre}
\begin{blocLivre}{ex-2-4}\begin{exercice}[1]{Ajouter 9, 19, 29 de tête}
Calculer mentalement $457+9$, $457+19$ et $262+29$.
\end{exercice}\end{blocLivre}
\begin{blocLivre}{ex-2-5}\begin{exercice}[1]{Estimer avant de calculer}
Sans poser l’opération, choisir le résultat de $4\,980+3\,120$ parmi : $810$ ; $8\,100$ ; $81\,000$. Justifier.
\end{exercice}\end{blocLivre}

\rubrique{Entraînement}
\begin{blocLivre}{ex-2-6}\begin{exercice}[2]{Regrouper pour aller vite}
Calculer $254+378+46$ en regroupant deux termes bien choisis.
\end{exercice}\end{blocLivre}
\begin{blocLivre}{ex-2-7}\begin{exercice}[2]{Une soustraction à trous}
Compléter : $\ldots-247=385$, puis $902-\ldots=345$.
\end{exercice}\end{blocLivre}
\begin{blocLivre}{ex-2-8}\begin{exercice}[2]{Le bon panneau}
Entre Villeneuve et Beaupré, un automobiliste lit « Beaupré 68 km » puis, plus loin, « Beaupré 29 km ». Quelle distance a-t-il parcourue entre les deux panneaux ?
\end{exercice}\end{blocLivre}
\begin{blocLivre}{ex-2-9}\begin{exercice}[2]{Avec un schéma en barres}
Un théâtre propose 850 places. 512 billets sont vendus le samedi et 218 le dimanche. Représenter la situation puis dire combien de places restent libres.
\label{fig-ex-2-9}\figureGeo{barres-9}
\end{exercice}\end{blocLivre}
\begin{blocLivre}{ex-2-10}\begin{exercice}[2]{Le compte est-il bon ?}
Issa calcule $4\,003-1\,587=3\,584$. Vérifier avec une addition, corriger s’il le faut.
\end{exercice}\end{blocLivre}

\rubrique{Problèmes}
\begin{blocLivre}{ex-2-11}\begin{exercice}[3]{La bibliothèque}
Une bibliothèque possède $12\,540$ livres. Elle en prête $3\,275$, puis reçoit un don de $1\,860$ livres. Combien de livres sont présents sur les étagères ?
\end{exercice}\end{blocLivre}
\begin{blocLivre}{ex-2-12}\begin{exercice}[3]{Qui a le plus ?}
Nina a 325 cartes. Elle a 68 cartes de plus que Marco. Combien Marco a-t-il de cartes ? Combien de cartes ont-ils ensemble ?
\label{fig-ex-2-12}\figureGeo{barres-12}
\end{exercice}\end{blocLivre}
\begin{blocLivre}{ex-2-13}\begin{exercice}[3]{La cagnotte de la classe}
Une classe doit réunir $1\,500$ euros pour un séjour. La vente de gâteaux rapporte 480 euros, la tombola 655 euros. Combien manque-t-il encore ?
\end{exercice}\end{blocLivre}
\begin{blocLivre}{ex-2-14}\begin{exercice}[3]{Trois étapes de course}
Un rallye de $2\,000$ km se court en trois étapes. La première mesure 712 km, la deuxième 56 km de moins que la première. Quelle est la longueur de la troisième étape ?
\end{exercice}\end{blocLivre}
\begin{blocLivre}{ex-2-15}\begin{exercice}[3]{Le vrai du faux}
Chloé affirme : « Pour calculer $757-398$, je calcule $759-400$, c’est plus simple. » A-t-elle raison ? Expliquer puis donner le résultat.
\end{exercice}\end{blocLivre}

\clearpage
